\documentclass[UTF8,a4paper,zihao=-4,AutoFakeBold]{ctexart}
\usepackage[a4paper,top=2.54cm,bottom=2.54cm,left=3.17cm,right=3.17cm]{geometry}
\linespread{1.55}
\usepackage{fontspec}
\usepackage{xeCJK}
\setCJKfamilyfont{fsgb}{fsGB2312.ttf}
\newcommand{\fsgb}{\CJKfamily{fsgb}}
\setCJKfamilyfont{fzxbsong}{fzxbsongti.TTF}
\newcommand{\fzxbsong}{\CJKfamily{fzxbsong}}
\ctexset{autoindent=2em}
\usepackage{amsmath}
\usepackage{amssymb}
\usepackage{graphicx}
\usepackage{float}
\usepackage[section]{placeins}
\usepackage{needspace}
\usepackage{caption}
\usepackage{subcaption}
\usepackage{booktabs}
\usepackage{array}
\usepackage{multirow}
\usepackage{tabularx}
\usepackage{longtable}
\usepackage{threeparttable}
\captionsetup{font={small,singlespacing},labelsep=quad}
\captionsetup[table]{font={small,singlespacing},justification=centering}
\captionsetup[figure]{font={small,singlespacing},justification=centering}
\newcolumntype{C}[1]{>{\zihao{5}\songti\centering\arraybackslash}p{#1}}

\usepackage{titlesec}
\usepackage{pifont}
\ctexset{
    section={format=\heiti\zihao{-3},name={,、},aftername={},number=\chinese{section},beforeskip=0pt,afterskip=0pt},
    subsection={format=\kaishu\zihao{4},name={(,)},aftername={},number=\chinese{subsection},beforeskip=0pt,afterskip=12pt},
    subsubsection={format=\songti\zihao{-4}\bfseries,aftername={},number=\arabic{subsubsection}.,beforeskip=0pt,afterskip=6pt}
}
% Level-4 heading via \ssection command
\newcounter{ssection}[subsubsection]
\renewcommand{\thessection}{\ding{\numexpr171+\value{ssection}\relax}}
\newcommand{\ssection}[1]{%
  \refstepcounter{ssection}%
  \par%
  \noindent%
  {\thessection#1}%
  \par\noindent%
}
\titlespacing{\section}{0pt}{0pt}{0pt}
\titlespacing{\subsection}{2em}{0pt}{0pt}
\titlespacing{\subsubsection}{2em}{0pt}{0pt}
\setlength{\parindent}{2em}

\usepackage{fancyhdr}
\pagestyle{fancy}
\fancyhf{}
\renewcommand{\headrulewidth}{0pt}
\fancyfoot[C]{\thepage}
\fancypagestyle{plain}{\fancyhf{}\renewcommand{\headrulewidth}{0pt}\fancyfoot[C]{\thepage}}

\usepackage[backref]{hyperref}
\hypersetup{hidelinks}
\usepackage[numbers,square,super]{natbib}
\renewcommand{\refname}{\heiti\zihao{4}参考文献}
\renewcommand{\bibfont}{\songti\zihao{-4}}

\usepackage{setspace}
\usepackage{tikz}
\usetikzlibrary{arrows.meta,positioning,shapes.geometric,calc}
\usepackage{listings}
\lstset{language=Python,basicstyle=\small\ttfamily,breaklines=true,frame=single,numbers=left,numberstyle=\tiny}

\setcounter{tocdepth}{3}
\usepackage{titletoc}
\renewcommand{\contentsname}{\heiti\zihao{4}\centering 目\hspace{1em}录}
\titlecontents{section}[0pt]{\addvspace{0pt}\songti\zihao{-4}}{\thecontentslabel\hspace{0em}}{}{\titlerule*[0.5pc]{.}\contentspage}
\titlecontents{subsection}[2em]{\addvspace{0pt}\songti\zihao{-4}}{\thecontentslabel\hspace{0em}}{}{\titlerule*[0.5pc]{.}\contentspage}
\titlecontents{subsubsection}[4em]{\addvspace{0pt}\songti\zihao{-4}}{\thecontentslabel\hspace{0.4em}}{}{\titlerule*[0.5pc]{.}\contentspage}


\begin{document}

% ==================== 封面 ====================
\thispagestyle{empty}
{\heiti\zihao{3} 作品编号：TJJM}\par
\vspace*{5.5cm}
\begin{center}
\fzxbsong
{\zihao{-2} \textbf{[竞赛年份]}年（第[届数]届）全国大学生统计建模大赛}\\
\vspace*{20pt}
{\zihao{1}参~赛~作~品}
\end{center}\par
\vspace*{4.5cm}
\begin{table}[h]
\centering
\zihao{-3}
\fsgb
\renewcommand{\arraystretch}{2}
\setlength{\tabcolsep}{5mm}{
\begin{tabular}{rc}
参赛学校：&[学校名称]\\\cline{2-2}
论文题目：&[论文标题]\\\cline{2-2}
参赛队员：&[队员1]~[队员2]~[队员3]\\\cline{2-2}
指导老师：&[指导老师]\\\cline{2-2}
\end{tabular}}
\end{table}

% ==================== 标题页 ====================
\clearpage
\begin{center}
\zihao{3}\fzxbsong\setlength{\baselineskip}{24pt}
[论文标题]
\end{center}

% ==================== 中文摘要 ====================
\section*{\centering\heiti\zihao{4} 摘要}
\addcontentsline{toc}{section}{摘要}
\songti\zihao{-4}

[中文摘要内容：500-700字]

\heiti\zihao{-4}\setlength{\baselineskip}{24pt} 关键词：
\songti\zihao{-4}\setlength{\baselineskip}{24pt} [关键词1]；[关键词2]；[关键词3]；[关键词4]；[关键词5]

% ==================== 英文摘要 ====================
\clearpage
\section*{\centering\bfseries\zihao{4} Abstract}
\addcontentsline{toc}{section}{Abstract}
\zihao{-4}

[English abstract, 350-500 words]

\textbf{Keywords:} [keyword1]; [keyword2]; [keyword3]; [keyword4]; [keyword5]

% ==================== 目录 ====================
\clearpage
\tableofcontents


% ==================== 表格与插图清单 ====================
\clearpage
\section*{\centering\heiti\zihao{4} 表格与插图清单}
\addcontentsline{toc}{section}{表格与插图清单}
\listoftables
\listoffigures

% ==================== 正文 ====================
\clearpage
\begin{center}
\zihao{3}\fzxbsong\setlength{\baselineskip}{24pt}
[论文标题]
\end{center}

% Chapter files: Claude designs chapter structure based on TOPIC_PLAN.md.
% Rename/add/remove \input lines to match actual chapter files.
% Typical: 5-7 chapters (introduction + 3-5 middle chapters + conclusion).
\input{sections/1_introduction.tex}
\input{sections/2_method.tex}
\input{sections/3_data.tex}
\input{sections/4_analysis.tex}
\input{sections/5_results.tex}
\input{sections/6_conclusion.tex}

% ==================== 参考文献 ====================
\addcontentsline{toc}{section}{参考文献}
\bibliography{references}
\bibliographystyle{plainnat}

% ==================== 附录 ====================
\addcontentsline{toc}{section}{附录}
\section*{\centering\heiti\zihao{4} 附录}

[附录内容：核心代码、补充表格等]

% ==================== 致谢 ====================
\section*{\centering\heiti\zihao{4} 致谢}
\addcontentsline{toc}{section}{致谢}
\songti\zihao{-4}

[致谢内容]

\end{document}
